\chapter{Характерно-разрешённый радиационный форм-фактор компакттона}
\label{ch:v6-spectral-transition-discrete-compacton-character-resolved-radiation-form-factor-gate}

\section{Решающий вопрос}

Предыдущий гейт установил, что оператор перехода
\begin{equation}
 T=|+i\rangle\langle-i|
\end{equation}
несёт аффинный характер $-1$. Компактная граничная мода сама даёт только
когерентный обмен. Однако полный нелинейный дискретный родитель уже имеет
пространственный радиационный континуум. Здесь вычисляется не абстрактная
плотность состояний, а точный форм-фактор характерного направления в
полном вещественном флоке-якобиане четырёхшагового возврата.

Антикруговая сверка использует
\path{version6_spectral_transition_discrete_compacton_stability_quantization_gate},
\path{version6_spectral_transition_discrete_compacton_dynamical_capture_gate},
\path{version6_spectral_transition_compacton_c4_affine_selector_admissibility_gate},
\path{version6_spectral_transition_discrete_compacton_c4_boundary_eta_dissipation_gate}
и
\path{version6_spectral_transition_discrete_compacton_energy_degeneracy_boundary_overlap_gate}.

\section{Две вещественные квадратуры одного характерного перехода}

Пусть $\Psi_+$ и $\Psi_-$ --- ортонормированные точные компакттоны с
одношаговыми фазами $+i$ и $-i$. Для эрмитова смесителя
\begin{equation}
 H_\theta=e^{i\theta}T+e^{-i\theta}T^\dagger
\end{equation}
касательное направление из $\Psi_+$ равно
\begin{equation}
 q_\theta=-i e^{-i\theta}\Psi_-.
 \label{eq:v6-character-radiation-quadrature}
\end{equation}

Обозначим через $M=D(F^4)_{\Psi_+}$ вещественную производную
четырёхшагового возврата. Численный скан семнадцати фаз даёт точную
формулу
\begin{equation}
 \|P_{\rm out}Mq_\theta\|
 =2\pi|\sin\theta|
 \label{eq:v6-character-radiation-sine-law}
\end{equation}
с максимальным остатком $6.13\cdot10^{-10}$.

При $\theta=0$ направление $-i\Psi_-$ лежит в касательной к уже
доказанному точному фазовому многообразию компакттонов. Поэтому
\begin{equation}
 M(-i\Psi_-)=-i\Psi_-
\end{equation}
с остатком $3.14\cdot10^{-12}$ и радиации нет.

Редуцированная относительная эта-фаза $-i$ соответствует
$\theta=-\pi/2$ в принятой ориентации. Тогда
\begin{equation}
 q_{-\pi/2}=\Psi_-,
 \qquad
 M\Psi_- = \Psi_-+r_1,
 \qquad
 \|r_1\|=2\pi.
 \label{eq:v6-character-eta-radiation-packet}
\end{equation}
То есть эта-ориентированная квадратура действительно имеет ненулевой
форм-фактор в пространственный канал.

\begin{proposition}[Форм-фактор не нулевой]
Полный флоке-якобиан различает две вещественные квадратуры перехода.
Фазовая квадратура является точной нейтральной касательной, а
эта-ориентированная амплитудная квадратура испускает пакет нормы $2\pi$
за один четырёхшаговый возврат.
\end{proposition}

\section{Пакеты являются исходящими}

Положим
\begin{equation}
 r_n=M^n\Psi_- - M^{n-1}\Psi_-.
\end{equation}
На решётке из $256$ узлов первые восемь пакетов удовлетворяют
\begin{equation}
 \|r_n\|=2\pi,
 \qquad
 |\langle r_m,r_n\rangle|<2.02\cdot10^{-8}
 \quad(m\ne n).
 \label{eq:v6-character-outgoing-orthogonal-packets}
\end{equation}
Их носители каждый цикл уходят ещё на четыре узла от ядра. До
периодического возврата имеем
\begin{equation}
 \|M^n\Psi_-\|^2
 =1+4\pi^2n
 \label{eq:v6-character-secular-emission}
\end{equation}
с максимальным остатком $9.20\cdot10^{-7}$.

Это доказывает наличие исходящего канала, но одновременно показывает
важную структуру: проекция на исходное характерное направление остаётся
равной единице. $\Psi_-$ не исчезает, а действует как секулярный источник
новых пакетов.

\section{Спектральная плотность и точный коэффициент}

На пустом фоне четырёхшаговый сдвиг имеет множители
\begin{equation}
 e^{\pm4ik}.
\end{equation}
Резонанс с единичным множителем возникает при
$k=0,\pi/2,\pi,3\pi/2$. Для первого пакета квадраты фурье-форм-факторов
равны
\begin{equation}
 (0,8\pi^2,0,8\pi^2)
\end{equation}
с машинными остатками. С учётом якобиана дисперсии $|d(4k)/dk|=4$
спектральная плотность при резонансе равна
\begin{equation}
 \rho_{\rm rad}(0)
 =\frac1{2\pi}\sum_{4k=0}
 \frac{\|\widehat r_1(k)\|^2}{4}
 =2\pi.
 \label{eq:v6-character-radiation-density}
\end{equation}
Поэтому формальный коэффициент золотого правила за один четырёхшаговый
цикл равен
\begin{equation}
 \Gamma_{\rm cycle}
 =2\pi\rho_{\rm rad}(0)|\delta|^2
 =4\pi^2|\delta|^2.
 \label{eq:v6-character-golden-coefficient}
\end{equation}

Исключение открытых каналов и появление антиэрмитовой части эффективного
оператора являются стандартным содержанием фешбаховского и фановского
формализма \cite{character-radiation-feshbach1958,
character-radiation-fano1961}. Для дискретных квантовых блужданий
непрерывный спектр и теория рассеяния также требуют явного анализа
спектральной меры, а не одного счёта размерностей
\cite{character-radiation-bisio2021}.

\section{Нелинейная проверка знака эффекта}

Формула~\eqref{eq:v6-character-golden-coefficient} сама по себе ещё не
говорит, что затухает именно нежелательная ветвь. Поэтому запущена полная
нелинейная унитарная эволюция состояния
\begin{equation}
 \frac{\Psi_++\delta\Psi_-}
 {\|\Psi_++\delta\Psi_-\|}
\end{equation}
на $512$ узлах в течение двадцати четырёхшаговых циклов.

При $\delta=10^{-4}$ получено
\begin{equation}
 1-P_{\rm core}=7.90747\cdot10^{-6},
\end{equation}
тогда как слабая экспоненциальная оценка даёт
$7.89565\cdot10^{-6}$. Совпадение подтверждает коэффициент $4\pi^2$.

Но характерный вес $\Psi_-$ не убывает:
\begin{equation}
 w_-^{\rm initial}=9.9999999\cdot10^{-9},
 \qquad
 w_-^{\rm final}=1.0030771\cdot10^{-8}.
\end{equation}
Убывает вероятность всего компакттонного ядра. При увеличении
$\delta$ возникает уже известный конечнопороговый уход компакттона.
Ортогональная квадратура $-i\Psi_-$ при $\delta=10^{-3}$ сохраняет ядро
с точностью $5.8\cdot10^{-15}$.

\begin{theorem}[Радиация имеет неправильный динамический смысл]
Коэффициент-свободный пространственный форм-фактор существует и его
безразмерный коэффициент равен $4\pi^2$. Однако линейная характерная
амплитуда не истощается: она секулярно излучает ортогональные пакеты.
Полная нелинейная динамика истощает компакттонное ядро, а не очищает его
до одной ветви. Поэтому найденный канал является не амплитудным
затуханием, а радиационной неустойчивостью.
\end{theorem}

\section{Почему это не скрытый \texorpdfstring{$\gamma$}{gamma}}

Число $4\pi^2$ не вставлено: оно является произведением точной нормы
пакета и спектральной плотности сдвига. В этом смысле произвольный
коэффициент перед уже выбранным открытым каналом действительно заменён
вычислимым форм-фактором.

Но физический rate всё ещё не замкнут по двум причинам:
\begin{enumerate}
 \item $\delta$ является амплитудой возбуждения линии $-1$; её вакуумная
 нормировка и начальное значение не выведены;
 \item физическое время одного шага $\Delta t$ остаётся свободным.
\end{enumerate}
Главное же препятствие сильнее масштабного: даже при заданной $\delta$
эффект разрушает ядро и оставляет нейтральную квадратуру.

Полная KO6 Real-пара дополнительно сокращает чистую эта-фазу, поэтому
ориентацию одной редуцированной квадратуры нельзя автоматически считать
глобальным выбором физической ветви.

\begin{nogocriterion}[Излучение не равно захвату]
Нельзя отождествлять ненулевую спектральную плотность с работающей
диссипацией. Для захвата необходимо, чтобы уменьшалась проекция на
нежелательную характерную моду и чтобы все физические поперечные
квадратуры имели отрицательную действительную скорость. Здесь первое не
выполнено, а вторая квадратура точно нейтральна.
\end{nogocriterion}

\section{Вердикт}

Поиск форм-фактора дал первый аналитически нормированный открытый канал
компакттонной ветви:
\begin{equation}
 \boxed{\rho_{\rm rad}(0)=2\pi,
 \qquad \Gamma_{\rm cycle}=4\pi^2|\delta|^2.}
\end{equation}
Но это не механизм рождения или выбора материи. Источник не затухает,
ядро теряет норму, а ортогональная квадратура остаётся на непрерывном
точном многообразии компакттонов. Следовательно, комбинация
``аффинный характер $-1$ + эта-фаза + существующая пространственная
радиация'' не заменяет амплитудный jump.

Следующий гейт
\path{version6_spectral_transition_discrete_compacton_branch_status_freeze_gate}
должен зафиксировать компакттон как точное локализованное решение и
полезный спектральный стенд, но закрыть его как автономный механизм
захвата и рождения материи.

\begin{protocol}
\begin{enumerate}
 \item точные ветви $\pm i$: \textbf{подтверждены};
 \item форм-фактор эта-ориентированной квадратуры: \textbf{ненулевой};
 \item норма исходящего пакета: \textbf{$2\pi$};
 \item спектральная плотность на резонансе: \textbf{$2\pi$};
 \item коэффициент золотого правила: \textbf{$4\pi^2$};
 \item ортогональность последовательных пакетов: \textbf{да};
 \item проекция источника убывает: \textbf{нет};
 \item характерный вес $w_-$ убывает: \textbf{нет};
 \item компакттонное ядро истощается: \textbf{да};
 \item ортогональная квадратура: \textbf{точно нейтральна};
 \item нелинейная слабая проверка коэффициента: \textbf{пройдена};
 \item асимптотический захват: \textbf{нет};
 \item ручной $\gamma$ заменён работающим автономным rate:
 \textbf{нет};
 \item ветвь как механизм рождения: \textbf{к замораживанию}.
\end{enumerate}
\end{protocol}

Машинный сертификат сохранён в
\path{s2t_v6_spectral_transition_discrete_compacton_character_resolved_radiation_form_factor_gate_results.json}.

\begin{thebibliography}{9}
\bibitem{character-radiation-feshbach1958}
H.~Feshbach,
\emph{Unified Theory of Nuclear Reactions},
Annals of Physics 5 (1958), 357--390,
\path{doi:10.1016/0003-4916(58)90007-1}.

\bibitem{character-radiation-fano1961}
U.~Fano,
\emph{Effects of Configuration Interaction on Intensities and Phase Shifts},
Physical Review 124 (1961), 1866--1878,
\path{doi:10.1103/PhysRev.124.1866}.

\bibitem{character-radiation-bisio2021}
A.~Bisio, N.~Mosco, and P.~Perinotti,
\emph{Scattering and Perturbation Theory for Discrete-Time Dynamics},
Physical Review Letters 126 (2021), 250503,
\path{arXiv:1912.09768}.
\end{thebibliography}